
%%Einstellungen für das Dokument: 

\renewcommand{\Autor}{Name}
\renewcommand{\Vorname}{Vorname}
\renewcommand{\Nachname}{Nachname}
\renewcommand{\Matrikelnummer}{Matrikelnummer}
\renewcommand{\Datum}{Abgabedatum}

%%Titel der Arbeit:
\renewcommand{\Titel}{Spannendes Thema}
\renewcommand{\TitelEN}{Exciting Title}


%%Typ der Arbeit
\renewcommand{\Dokumentenart}{Bachelorarbeit}
\renewcommand{\Pruefungsart}{Bachelorprüfung}

%% Prüfer und Firma:
\renewcommand{\Erstpruefer}{Erstprüfer: Prof. Dr. Ing. }
\renewcommand{\Zweitpruefer}{Zweitprüfer: Prof. Dr. Ing.}
\renewcommand{\Betreuer}{Industrieller Betreuer: Name}
\renewcommand{\Firmenname}{Firma XY AG} 
\renewcommand{\Firma}{     
        in Zusammenarbeit mit:\\
        \Firmenname\\
        Abteilung XX\\
        Straße Nr.\\
        PLZ Ort}  



%% ABSTRACT PAGE


%% Keywords DE/ENG
\renewcommand{\Stichworte}{Test 1, Test 2}
\renewcommand{\StichworteEN}{Test 1, Test 2}

%%Zusammenfassung DE/ENG
\renewcommand{\Zusammenfassung}{noch nicht viel passiert...}
\renewcommand{\ZusammenfassungEN}{not much happened yet...}

%% Optional: KI-Verwendungshinweis
\defbibnote{ai_disclosure}{%
  \small\itshape
  Hinweis zur Verwendung von KI-Werkzeugen: Bei der Anfertigung der vorliegenden Arbeit wurde das Large Language Model Claude unterstützend für das stilistische Lektorat, die Code-Generierung sowie die vorbereitende Literaturrecherche eingesetzt. Die kritische Evaluierung der Ergebnisse, die inhaltliche Synthese sowie die vollumfängliche wissenschaftliche Verantwortung verbleiben uneingeschränkt beim Autor.
}